\scanpage{87}
\margin{Catégories $\underline{\Gamma}$ et $\underleftarrow{\mathrm{Lim}}$}

\diagram{\begin{tikzcd}[ampersand replacement=\&,row sep=tiny]
\unsure{E|\mathcal E}\arrow[d]\\ B
\end{tikzcd}}

Catégorie fibrée

\diagram{\begin{tikzcd}[ampersand replacement=\&,row sep=large,column sep=large]
\&\underleftarrow{\mathrm{Lim}}\ E/B\arrow[r,hook,"i"]\arrow[d,"F_c^*"]\&\underline{\Gamma}(E/B)\arrow[d,"F^*"]\\
\underline{X'\in\mathrm{Ob}}\&\underleftarrow{\mathrm{Lim}}\ E'/B'\arrow[r,hook,no head,"i'"]\&\underline{\Gamma}(E'/B')
\end{tikzcd}}

\[
\begin{aligned}
\mathrm{Hom}(X',F_c^*(Y))&=\mathrm{Hom}(i'(X'),F^*(i(Y)))\\
&=\mathrm{Hom}(F_*\,i'(X'),i(Y))\\
&=\mathrm{Hom}(\underbrace{i_*F_*\,i'(X')}_{F_*^c(X')},Y)
\end{aligned}
\]

\textcircled{1} \uline{Si $\exists$ adjoint à g. $i_{\unsure{+|*}}$ de $i$}, et si $\exists$ $F_{\unsure{*}}$ adjoint
à g. de $F^*$ (p.ex. si $E$ est \uline{aussi} cofibrée
sur $B$, et les ``petites'' $\unsure{\underrightarrow{\lim}|\underline{\lim}|\underleftarrow{\lim}}$ existent dans les fibres)
alors $\exists$ adjoint à gauche $F_*^c$ de ${F}_c^*$.

\[
\begin{aligned}
\mathrm{Hom}(F_c^*(Y),X')&=\mathrm{Hom}(F^*(i(Y)),i'(X'))\\
&=\mathrm{Hom}(i(Y),F_+\,i'(X'))\\
&=\mathrm{Hom}(Y,\underbrace{i_+F_+\,i'}_{F_+^c}(X'))
\end{aligned}
\]

\textcircled{2} \uline{Si} $\exists$ \uline{adjoint à dr.} $i_+$ \uline{de} $i$, et si $\exists$ adj. à dr.
$F_+$ de $F^*$ (p.ex. \deleted{\unsure{les fibres de $E/B$}} si les petites $\underleftarrow{\lim}$ existent ds

% archive sheet number: 86
